\documentclass[10pt,a4paper]{amsart}
\usepackage[margin=19mm]{geometry}
\usepackage{amsmath,amssymb,mathtools}
\usepackage[scr=rsfs,cal=euler]{mathalfa}
\usepackage{xfrac}
\usepackage[unicode,hidelinks,bookmarks=false]{hyperref}
\newcommand{\ie}{i.e.\ }
\newcommand{\cf}{cf.\ }
\newcommand{\resp}{respectively\ }
\newcommand{\Subsection}{\subsection}
\providecommand{\nobreakdash}{\nobreak\hbox{-}}
\setlength{\emergencystretch}{2em}
\multlinegap0pt
\usepackage{graphics}
\usepackage{epsfig}
\usepackage{amscd}
\numberwithin{equation}{section}
\usepackage{mathrsfs}
\usepackage{color}
\usepackage{bm}
\usepackage{float}
\usepackage{tikz}
\usepackage{multirow}
\usepackage{verbatim}
\usepackage{placeins}
\usetikzlibrary{matrix,arrows,arrows.meta,calc}
\usetikzlibrary{decorations.pathreplacing,decorations.markings}
\usetikzlibrary{shapes,positioning}
\tikzstyle arrowstyle=[scale=1]
\tikzstyle directed=[postaction={decorate,decoration={markings,
    mark=at position .5 with {\arrow[arrowstyle]{stealth}}}}]
\tikzstyle reverse directed=[postaction={decorate,decoration={markings,
    mark=at position .65 with {\arrowreversed[arrowstyle]{stealth};}}}]
\newtheorem{thm}{Theorem}[section]
\newtheorem{cor}[thm]{Corollary}
\newtheorem{prop}[thm]{Proposition}
\newtheorem{conj}[thm]{Conjecture}
\newtheorem{prob}[thm]{Problem}
\newtheorem{lem}[thm]{Lemma}
\newtheorem{rem}[thm]{Remark}
\newtheorem{defn}[thm]{Definition}
\newtheorem{ex}[thm]{Example}
\numberwithin{equation}{section}
\newcommand{\C}{\mathbb{C}}
\newcommand{\Q}{\mathbb{Q}}
\newcommand{\R}{\mathbb{R}}
\newcommand{\Z}{\mathbb{Z}}
\newcommand{\Frac}{\mathrm{Frac}}
\newcommand{\A}{\mathcal{A}}
\newcommand{\B}{\mathcal{B}}
\newcommand{\CC}{\mathcal{C}}
\newcommand{\F}{\mathcal{F}}
\newcommand{\FF}{\mathscr{F}}
\newcommand{\E}{\mathscr{E}}
\newcommand{\Flag}{\mathrm{Flag}}
\newcommand{\T}{\mathcal{T}}
\newcommand{\Ch}{\mathrm{Ch}}
\newcommand{\VG}{\mathrm{VG}}
\newcommand{\OS}{\mathrm{OS}}
\newcommand{\sVG}{\mathscr{VG}}
\newcommand{\sOS}{\mathscr{OS}}
\newcommand{\Mag}{\mathrm{Mag}}
\newcommand{\geod}{\mathrm{geod}}
\newcommand{\gs}{\mathrm{gs}}
\newcommand{\qq}{\mathbf{q}}
\newcommand{\cro}{\mathrm{cr}}
\newcommand{\La}{\mathscr{L}}
\newcommand{\PP}{\mathcal{P}}
\newcommand{\M}{\mathcal{M}}
\newcommand{\LL}{\mathcal{L}}
\newcommand{\id}{\text{id}}
\newcommand{\supp}{\mathrm{supp}}
\DeclareMathOperator{\rank}{rank}
\DeclareMathOperator{\diam}{diam}
\DeclareMathOperator{\codim}{codim}
\DeclareMathOperator{\colim}{colim}
\DeclareMathOperator{\coker}{coker}
\DeclareMathOperator{\sgn}{sgn}
\DeclareMathOperator{\gr}{gr}
\DeclareMathOperator{\Hilb}{Hilb}
\DeclareMathOperator{\NBC}{NBC}
\newcommand{\Mod}{\mathbf{Mod}}
\begin{document}
\title[Orlik--Solomon sheaf homology of geometric lattices]{Orlik--Solomon sheaf homology of geometric lattices}
\author{Secondary 05B35 ạte; Ye Liu}
\date{}
\maketitle
\noindent\emph{Source and licence.} Orlik--Solomon sheaf homology of geometric lattices. Version arXiv:2607.08494v1 (CC BY 4.0). This material is distributed under the Creative Commons Attribution 4.0 International licence, http://creativecommons.org/licenses/by/4.0/, as recorded in the arXiv OAI licence metadata for this exact version and shown on the abstract page https://arxiv.org/abs/2607.08494v1. Full rights notice: licenses/arxiv-2607.08494-CC-BY-4.0.txt.\par\smallskip
\noindent\textit{Editorial scope.} Counted unit: the original \texttt{lem} environment \texttt{-} at source lines 441--450 together with its complete proof at lines 451--453; the delivered document keeps source lines 264--454 so that every referenced label and every citation resolves inside it. Native editable source is the paper's own concatenated TeX; the AMS wrapper is editorial formatting only. No publisher class, upstream script or unrelated artwork is redistributed.
\par\smallskip\noindent\textit{Reference note.} This article ships no compiled bibliography (biblatex); the entries delivered here are composed from the author's own BibTeX (.bib) fields, with no field invented and no entry dropped.
\begin{prop}[\cite{Gabriel1967} Appendix II, see also \cite{Everitt2022} Section 2.2]\label{chaincomplex}
    We have an isomorphism
\[
H_*(P;\FF)\cong H(T_*(P;\FF)).
\]
\end{prop}

Some easy observations of poset sheaf homology are as follows.
\begin{lem}[\cite{Everitt2022} Lemma 1]\label{basiclemma}
    Let $P$ be a poset and $\FF$ a sheaf on $P$.
    \begin{enumerate}
        \item If $P$ is graded, then $H_i(P;\FF)\neq 0$ only if $0\leq i\leq \rank P$.
        \item If $P$ has a top or a bottom, and $\FF=\underline{M}$ is the constant sheaf of $R$-module $M$, then $H_0(P;\underline{M})=M$ and $H_i(P;\underline{M})=0$ for $i>0$.
        \item If $P$ has a bottom $\hat{0}$, then $H_0(P;\FF)=\FF(\hat{0})$ and $H_i(P;\FF)=0$ for $i>0$.
        \item If $P$ has a top $\hat{1}$, and $\FF(\hat{1})=0$, then $H_*(P;\FF)\cong H_*(P\setminus\{\hat{1}\};\FF)$.
    \end{enumerate}
\end{lem}

Let $P$ be a poset with a bottom $\hat{0}$ and $\FF$ a sheaf on $P$. We define an augmentation map 
\[
\epsilon: T_0(P\setminus\{\hat{0}\};\FF)\to\FF(\hat{0})
\]
as the sum of structure maps $\FF^x_{\hat{0}}$ over $x\in P\setminus\{\hat{0}\}$. Then we obtain the augmented complex $\widetilde{T}_*(P\setminus\{\hat{0}\};\FF)$, whose homology $\widetilde{H}_*(P\setminus\{\hat{0}\};\FF)$ is called the reduced homology of $P\setminus\{\hat{0}\}$ with coefficients in the sheaf $\FF$. The map $\epsilon$ induces $\epsilon_*: H_0(P\setminus\{\hat{0}\};\FF)\to \FF(\hat{0})$, which coincides with the map $\colim^{P\setminus\{\hat{0}\}}\FF\to\FF(\hat{0})$ induced by $\FF^x_{\hat{0}}$ using the universality of colimit. We have
\[
\widetilde{H}_i(P\setminus\{\hat{0}\};\FF)\cong \begin{cases}
    H_i(P\setminus\{\hat{0}\};\FF), & i>0\\
    \ker\epsilon_*, & i=0\\
    \coker\epsilon_*, & i=-1
\end{cases}
\]


\subsection{OS algebra and OS sheaf}

Let $L$ be a geometric lattice. Fix a total order of the atom set $A=L_1=\{a_1,\ldots,a_n\}$ and introduce a set of symbols $E=\{e_1,\ldots,e_n\}$ in one-to-one correspondence to $A$. 
Let $R$ be a ring and $\wedge^\bullet_RE$ be the exterior algebra over $R$ generated by degree $1$ generators $\{e_1,\ldots,e_n\}$. For an index set $I=\{i_1,\ldots,i_k\}\subseteq [n]=\{1,2,\ldots,n\}$ with $i_1<\cdots<i_k$, write $A_I=\{a_{i_1},\ldots,a_{i_k}\}$, $E_I=\{e_{i_1},\ldots,e_{i_k}\}$ and $e_I=e_{i_1}\cdots e_{i_k}\in \wedge^k_RE$.
The degree $k$ piece $\wedge_R^kE$ is a free $R$-module with basis $\{e_I\mid I\subseteq [n],|I|=k\}$. There is a boundary operator $\partial: \wedge^k_RE\to \wedge^{k-1}_RE$ defined by
\[
\partial(e_{i_1}\cdots e_{i_k})=\sum_{j=1}^k(-1)^{j-1}e_{i_1}\cdots \widehat{e_{i_j}}\cdots e_{i_k},
\]
where $\widehat{e_{i_j}}$ means deleting $e_{i_j}$. 
We say $A_I(I\subseteq [n])$ is a {\it circuit} if it is minimally dependent.
\begin{defn}
Let $L$ be a geometric lattice with atom set $A$. Define the Orlik--Solomon (OS) algebra of $L$ by
\[
\OS^\bullet_R(L):=\wedge^\bullet_RE/\langle \partial e_I\mid A_I \text{ is a circuit}\rangle.
\]
Note that the ideal in the quotient is graded and thus $\OS^\bullet_R(L)$ is graded.
\end{defn}

Orlik--Solomon \cite{Orlik1980} proved that when $L=L(\A)$ is the intersection lattice of a complex hyperplane arrangement $\A$, the OS algebra is isomorphic to the cohomology ring of the complement $M(\A)$. In particular, the Hilbert series of the OS algebra is equal to the Poincar\'e polynomial of the complement.

Let $L$ be a general geometric lattice. For a circuit $A_I=\{a_{i_1},\ldots,a_{i_k}\}(i_1<\cdots<i_k)$ of $L$, its broken circuit is $A_{I\setminus\{\min I\}}=\{a_{i_2},\ldots,a_{i_k}\}$. A subset $A_I$ is said to be NBC, or no-broken-circuit, if it contains no broken circuit. Let us denote by $\NBC(L)$ the set of NBC subsets of $A$, and $\NBC^p(L)=\{A_I\in\NBC(L)\mid |I|=p\}$.

\begin{prop}[\cite{Orlik1992} Chapter 3]\label{OSproperties}
Let $L$ be a geometric lattice.
\begin{enumerate}
    \item The graded piece $\OS^p_R(L)$ is a free $R$-module with basis $\{\overline{e_I}\mid A_I\in \NBC^p(L)\}$.
    \item The Hilbert series of $\OS^\bullet_R(L)$ is equal to the Poincar\'e polynomial of $L$: $\Hilb(\OS^\bullet_R(L),t)=\pi(L,t)$.
    \item The natural map $\bigoplus_{x\in L_p}\OS^p_R(L_{\leq x})\to\OS^p_R(L)$ is an isomorphism of $R$-modules.
\end{enumerate}

\end{prop}


Our goal is to define a sheaf on $L$ using OS algebras. Let $x\in L$ be a flat, then $L_{\leq x}$ is itself a geometric lattice with atom set $A_x:=A(L_{\leq x})=A(L)\cap L_{\leq x}$. Let us assign the OS algebra $\OS_R^\bullet(L_{\leq x})$ to $x$. We denote by $E_x\subseteq E$ the set of generators of $\OS_R^\bullet(L_{\leq x})$. When $x\leq y$ in $L$, we have $L_{\leq x}\subseteq L_{\leq y}$, $A_x\subseteq A_y$ and $E_x\subseteq E_y$. The inclusion $E_x\hookrightarrow E_y$ induces an inclusion of exterior algebras $\wedge^\bullet_RE_x\hookrightarrow\wedge^\bullet_RE_y$ preserving dependency, hence induces an inclusion of OS algebras $\iota_{x,y}:\OS_R^
\bullet(L_{\leq x})\hookrightarrow \OS_R^
\bullet(L_{\leq y})$. Conversely, the map $\pi_{y,x}:E_y\to E_x$
    \[
    \pi_{y,x}(e_i)=\begin{cases}
        e_i, & e_i\in E_x,\\
        0, & e_i\in E_y\setminus E_x,
    \end{cases}
    \]
induces a well-defined surjective graded algebra homomorphism
\[
    \pi_{y,x}:\OS_R^\bullet(L_{\leq y})\to \OS_R^\bullet(L_{\leq x}).
    \]
In fact, let $A_I$ be a circuit in $L_{\leq y}$. If $A_I\subseteq A_x$, then $\partial e_I\in \wedge^\bullet_RE_y$ is mapped to the same element in $\wedge^\bullet_RE_x$. If $A_I\not\subseteq A_x$, then $|A_I\setminus A_x|\geq 2$ since if $A_I\setminus A_x=\{a\}$, the condition $A_I$ is a circuit would force $a\in A_x$. Therefore every monomial occurring in $\partial e_I$ contains at least one generator $e_i\in E_y\setminus E_x$ and $\partial e_I$ is mapped to $0$. 

\begin{defn}
    Let $L$ be a geometric lattice. We define a sheaf $\sOS_R^{\bullet}$ of graded $R$-algebras on $L$ by
    \[
    \sOS_R^{\bullet}(x):=\OS_R^\bullet(L_{\leq x}), ~x\in L,
    \]
    and the structure map
    \[
    (\sOS_R^\bullet)^y_x=\pi_{y,x}: \sOS_R^\bullet(y)=\OS_R^\bullet(L_{\leq y})\to\OS_R^\bullet(L_{\leq x})=\sOS_R^\bullet(x), ~x,y\in L, x\leq y.
    \]
    We call $\sOS_R^\bullet$ the Orlik--Solomon sheaf on $L$. The degree $p$ part is the sheaf of $R$-modules
    \[
    \sOS_R^p(x)=\OS_R^p(L_{\leq x}),
    \]
    with structure maps restricted to degree $p$ part.
\end{defn}


\section{Results}
Let $L$ be a geometric lattice of rank $\ell\geq 2$ and $P:=L\setminus\{\hat{0}\}$. The main result of this paper is the computation of the sheaf homology of $P$ with coefficients in the OS sheaf. For ease of notation, we shall fix $R=\Z$ and omit the reference to $R$ in the subscript.


\subsection{Statement of the result}


We first have some immediate observations. 
If $p=0$, $\sOS^0$ is the constant sheaf $\underline{\Z}$, then $\widetilde{H}_*(P;\sOS^0)=0$ by Lemma \ref{basiclemma} (2). 
If $p>0$, $\sOS^p(\hat{0})=0$, the augmentation map $\epsilon: S_0(P;\sOS^p)\to \sOS^p(\hat{0})$ is trivial and hence $H_*(P;\sOS^p)=\widetilde{H}_*(P;\sOS^p)$. 
Therefore we shall compute $H_i(P;\sOS^p)$ for $p>0$.

\begin{thm}\label{main}
    Let $p>0$. Then
    \[
    H_i(P;\sOS^p)\cong \begin{cases}
        \bigoplus_{z\in L_p}\OS^p(L_{\leq z})\otimes_{\Z}\widetilde{H}_{\ell-2}(\Delta Q_z;\Z),~&i=\ell-1,\\
        0,& i\neq \ell-1,
    \end{cases}
    \]
    where $Q_z=\{x\in P\mid z\not\leq x\}$. Furthermore $\OS^p(L_{\leq z})$ is free abelian of rank $(-1)^{\rank z}\mu(\hat{0},z)$ and $\widetilde{H}_{\ell-2}(\Delta Q_z;\Z)$ is free abelian of rank $b_z=(-1)^{\ell-2}\sum_{x\in L_{\geq z}}\mu(\hat{0},x)$. Therefore $H_{\ell-1}(P;\sOS^p)$ is free abelian of rank $\sum_{z\in L_p}(-1)^{\rank z}\mu(\hat{0},z)b_z$.
\end{thm}

Define the graded OS sheaf Euler characteristic of $L$ by
\[
\chi^\OS(L,t):=\sum_{p}\chi^p(L)t^p,~\chi^p(L):=\sum_i(-1)^i\rank_{\Z}H_i(P;\sOS^p).
\]
For $x\in L$, recall the Poincar\'e polynomial of $L_{\leq x}$
\[
\pi_{L_{\leq x}}(t)=\sum_{z\in L_{\leq x}}\mu(\hat{0},z)(-t)^{\rank z}.
\]
\begin{cor}
    We have
    \[
    \chi^\OS(L,t)=1-\sum_{x\in L}\mu(\hat{0},x)\pi_{L_{\leq x}}(t).
    \]
\end{cor}
\begin{proof}
    By Theorem \ref{main},
    \[
    \chi^p(L)=\begin{cases}
        -\sum_{z\in L_p}(-1)^p\mu(\hat 0,z)\sum_{x\in L_{\geq z}}\mu(\hat0,x), & p>0;\\
        1, & p=0.
    \end{cases}
    \]
    Then
    \begin{align*}
        \chi^\OS(L,t)&=1-\sum_{p\geq 1}\left(\sum_{z\in L_p}(-1)^p\mu(\hat 0,z)\sum_{x\in L_{\geq z}}\mu(\hat0,x)\right)t^p\\
        &=1-\sum_{x\in L\setminus\{\hat0\}}\mu(\hat0,x)\left(\sum_{z\in (\hat0,x]}\mu(\hat0 ,z)(-t)^{\rank z}\right)\\
        &=1-\sum_{x\in L\setminus\{\hat0\}}\mu(\hat0,x)(\pi_{L_{\leq x}}(t)-1)\\
        &=1-\sum_{x\in L\setminus\{\hat0\}}\mu(\hat0,x)\pi_{L_{\leq x}}(t)+\sum_{x\in L\setminus\{\hat0\}}\mu(\hat0,x)
    \end{align*}
    Since $\sum_{x\in L}\mu(\hat0,x)=0,\mu(\hat 0,\hat0)=1$ and $\pi_{\hat0}(t)=1$, the desired formula follows.
\end{proof}

\subsection{Proof of the main theorem}
Let $z\in L_p$, define its top local OS group
\[
B_z:=\OS^p(L_{\leq z}).
\]
It is free abelian of rank $\rank_{\Z}B_z=(-1)^{p}\mu(\hat{0},z)$ (Proposition \ref{OSproperties}(1)(2)). With the total order of atoms, $B_z$ has the usual NBC basis indexed by the $p$-element NBC sets whose join is $z$. The Brieskorn decomposition (Proposition \ref{OSproperties}(3)) gives, for every $x\in L$,
\[
\sOS^p(x)\cong \bigoplus_{z\in L_p\cap L_{\leq x}}B_z.
\]
Under this decomposition, $\pi_{y,x}:\sOS^p(y)\to\sOS^p(x)$ is the projection onto those summands with $z\leq x$. This suggests a decomposition of the degree $p$ OS sheaf
\[
\sOS^p\cong \bigoplus_{z\in L_p} B_z\otimes_{\Z}\E_z,
\]
where $\E_z$ is the indicator sheaf of $L_{\geq z}$,
\[
\E_z(x)=\begin{cases}
    \Z, & z\leq x,\\
    0, & z\not\leq x,
\end{cases}
\]
with identity maps within $L_{\geq z}$ and zero maps otherwise. Then we have for $p>0$,
\[
H_i(P;\sOS^p)\cong \bigoplus_{z\in L_p} H_i(P;B_z\otimes_{\Z}\E_z)\cong \bigoplus_{z\in L_p}B_z\otimes_{\Z}H_i(P;\E_z).
\]
The computation is reduced to $H_i(P;\E_z)$. We use Proposition \ref{chaincomplex}.

\begin{lem}
    There is a natural chain isomorphism
    \[
    T_*(P;\E_z)\cong C_*(\Delta P,\Delta Q_z;\Z).
    \]
    Consequently,
    \[
    H_i(P;\E_z)\cong \widetilde{H}_{i-1}(\Delta Q_z;\Z).
    \]
\end{lem}

\begin{proof}
    For a chain $\sigma=(x_n<\cdots< x_0)$, the contribution to the complex $T_*$ is $\E_z(x_0)$, which is nontrivial ($\Z$) precisely when $x_0\in P_{\geq z}$. Since $P_{\geq z}$ is a principal filter, this is equivalent to the chain having at least one vertex in $P_{\geq z}$, or equivalently to $\sigma$ being a simplex of $\Delta P$ not lying in $\Delta Q_z$, where $Q_z=P\setminus P_{\geq z}$. Hence the chain groups are isomorphic. We also need to verify the boundary operators agree. If deleting the top vertex $x_0$ leaves a new top $x_1$ in $P_{\geq z}$, the sheaf map is the identity. If the new top $x_1$ is in $Q_z$, the sheaf map is zero, exactly as in the relative quotient. Since $P$ has maximum $\hat{1}$, $\Delta P$ is contractible. The long exact sequence then gives the result.
\end{proof}


\begin{thebibliography}{99}
\bibitem[]{Everitt2022} Everitt, Brent; Turner, Paul. Deletion-restriction for sheaf homology of graded atomic lattices. Adv. Math. 402 (2022) Paper No. 108354, 24
\bibitem[]{Gabriel1967} Gabriel, P.; Zisman, M.. Calculus of fractions and homotopy theory. Band 35 (1967) x+168
\bibitem[]{Orlik1980} Orlik, Peter; Solomon, Louis. Combinatorics and topology of complements of hyperplanes. Invent. Math. 56 (1980) 167--189
\bibitem[]{Orlik1992} Orlik, Peter; Terao, Hiroaki. Arrangements of hyperplanes. 300 (1992) xviii+325
\end{thebibliography}
\end{document}
